\subsection{Inducing the shift}\label{sec:kinduced}
The contrast between the two datasets is consistent with the novel-attack explanation but does not test
it, because the datasets differ in more than their splits. We therefore induce the shift on UNSW-NB15
itself: whole attack categories are removed from the 2{,}000-row training sample and left in the
20{,}000-row test sample, and the protocol of Table~\ref{tab:twosplits} is repeated. Two holdouts guard
against a lucky choice. Holdout A removes Exploits, Fuzzers and Reconnaissance, which then make up $25\%$
of the test rows; holdout B removes Generic, DoS, Analysis and Backdoor, $29\%$ of the test rows.
Table~\ref{tab:induced} gives the result.

Unseen attack types are necessary for the failure but not sufficient. Under holdout A standard CV still
ranks the grid correctly for every kernel ($r=+0.85$ to $+0.99$) and its choices lie within $0.009$
ROC-AUC of the best setting in the grid. Under holdout B the NSL-KDD picture returns. Every kernel loses
between $0.06$ and $0.16$ ROC-AUC at its best setting, the correlation between CV score and test ROC-AUC
falls to $+0.62$ to $+0.69$, and the loss from selection is unequal across kernels: the CV-selected
raw-feature RBF kernel sits $0.108$ below its oracle ($0.794$ against $0.902$), the quantum kernel
$0.028$ below ($0.761$ against $0.789$). A comparison made at the CV-selected settings therefore
understates the classical kernel's lead by $0.08$ ROC-AUC, and at $1\%$ false positives it reverses the
order outright: the quantum kernel detects $0.302$ of attacks against $0.263$ for the RBF kernel, whose
best setting in the grid detects $0.465$. This is the mechanism of Section~\ref{sec:kcv} reproduced on
demand, on a dataset on which it was absent, by changing nothing but which attack types the training
sample contains. What separates the two holdouts is the bandwidth landscape: under holdout B the RBF
kernel's test ROC-AUC peaks sharply at $\gamma=0.3$ ($0.902$, against $0.794$ at $\gamma=0.1$ and
$0.864$ at $\gamma=1$) while its CV score is flat across those settings ($0.966$ to $0.967$), so the
selection criterion cannot resolve the one setting that generalises to the unseen categories.

Shift-aware selection does not repair holdout B. Leave-attack-types-out CV, which restored selection on
NSL-KDD, correlates with test ROC-AUC at $r=+0.43$ for the RBF kernel here and chooses a setting worse
than standard CV ($0.751$). Its folds can only hold out attack types the training sample contains
(Exploits, Fuzzers, Reconnaissance, Shellcode, Worms), and a shift towards those types is not the shift
towards Generic and DoS that the test set presents. On NSL-KDD the seventeen unseen types are variants
within the same four attack families as the seen ones, which is why holding out types reproduces the
official split there. Nor is the result an artefact of how the folds are built. The construction used
throughout assigns attack types to folds with GroupKFold, which leaves three usable folds on NSL-KDD and
one on UNSW-NB15; an alternative that splits the benign rows evenly and assigns attack types largest-first
to the emptiest fold makes all five folds usable on every dataset, but gives $r=+0.54$ instead of $+0.79$
for the RBF kernel on NSL-KDD, $r=-0.38$ under holdout B, and on the official UNSW-NB15 split, which has
no unseen types, it selects an over-smooth bandwidth for the RBF kernel (test ROC-AUC $0.859$ against
$0.938$ under standard CV). Shift-aware selection is a choice of target rather than a free improvement:
it buys robustness to the kinds of novelty the training data can imitate, at a cost when there is no
novelty to be robust to, and it cannot anticipate novelty of a kind the training data has never seen.
The general lesson is procedural. A kernel comparison under distribution shift should be reported under
more than one selection rule, with the oracle bound alongside, so that a reader can see how much of a
claimed difference is selection.

@@TAB_INDUCED@@
